\documentclass[11pt]{article}
\usepackage[margin=1in]{geometry}
\usepackage{amsmath,amssymb,amsthm,mathtools,booktabs,microtype}
\usepackage[colorlinks=true,linkcolor=blue,urlcolor=blue,citecolor=blue]{hyperref}
\hypersetup{pdftitle={Zeros of a whole family: a guide to the 7/8 argument},
 pdfsubject={Concrete averages in the 7/8 zero-free half-plane argument}}
\allowdisplaybreaks[2]
\setlength{\emergencystretch}{2em}
\newcommand{\OO}{\mathbb Z[\omega]}
\newcommand{\N}{\mathrm N}
\newcommand{\ch}{\chi_{-3}}
\newcommand{\ind}{\mathbf 1}
\newcommand{\starred}{\sideset{}{^*}\sum}
\newcommand{\sixsum}{\sideset{}{^{(6)}}\sum}
\newtheorem{theorem}{Theorem}
\newtheorem{proposition}{Proposition}
\newtheorem{lemma}{Lemma}
\title{Zeros of a whole family:\\an accessible guide to the $7/8$ argument for $\zeta(s)$}
\author{Claude Opus 5.5 \quad and \quad OpenAI Astra\\[1ex]
{\normalsize at the behest of Michael O.\ Rubinstein}}
\date{}
\begin{document}
\maketitle
\begin{center}\setlength{\fboxsep}{6pt}
\fbox{\parbox{0.9\linewidth}{\small\textbf{A note from Michael Rubinstein.} These are working notes,
generated by Claude and Codex to help me read the papers on the quasi-Riemann hypothesis that
OpenAI released.}}
\end{center}

\begin{abstract}
We describe, in classical terms and with the averages written out, the
argument of \cite{OAI78} showing that $\zeta(s)$ has no zeros in
$\Re s>7/8$. Unlike the $11/12$ argument of the companion guide
\cite{guide1112}, it is a contradiction argument about the rightmost zero
of a whole family of $L$-functions over $\mathbb Q(\sqrt{-3})$. A single
average (the ``probe'') is evaluated in two exact ways. The $7/8$ step
makes the dual sum longer; short sums over primes, arranged so that a scalar
contribution cancels, make the longer sum affordable, and the rare bad rows are
counted with two different moment bounds.
\end{abstract}

\paragraph{The proof in six steps.}
\begin{enumerate}
\item Let $\beta_*$ be the supremum of the real parts of the zeros (in
 $1/2\le\Re s\le1$) of all $L$-functions $L(s,\eta)$ in a fixed family over the
 Eisenstein integers; the family contains $\zeta(s)L(s,\chi_{-3})$. Assume
 $\beta_*>7/8$ and aim for a contradiction.
\item For every member $\eta$, one builds an explicit average $I(Z)$, the
 \emph{probe}, of cubic Gauss sums against sextic residue symbols and
 Gauss sums, and evaluates it in two exact ways.
\item The first way (Cauchy--Schwarz; a planar large sieve for one factor and,
 for the other, the automorphy of the cubic theta function followed by a
 quadratic large sieve) shows that the probe is small: at most
 $Z^{3/16+\varepsilon}$. Small averaging scales help here, but they make the
 dual sum in the second way long.
\item The second way (Poisson summation) writes the probe as a sum over
 ``rows'' $u$; the row $u=1$ contains $1/L(s,\eta)$ and would be of size
 $Z^{\beta_*-11/16}$ if $L(s,\eta)$ had a zero near $\beta_*$, while every
 other row involves twisted $L$-functions, and a row can be large only if one
 of \emph{them} has a zero of bounded height well to the right of the critical line.
\item A zero forces two Dirichlet polynomials attached to that row (one with
 M\"obius coefficients, one without) to be unusually large, and two moment
 estimates show that such rows are rare. Short sums over primes built into
 the probe, arranged so that a row-independent term cancels, are what make the
 long dual sum affordable: they have full size on the row $u=1$, but on the other
 rows they are only square-root-size prime polynomials, which also sharpen the count.
\item Comparing the two evaluations shows that $1/L(s,\eta)$ continues a
 fixed distance to the left of $\beta_*$ for every $\eta$, contradicting
 the definition of $\beta_*$; hence $\beta_*\le7/8$, and in particular
 $\zeta(s)\ne0$ for $\Re s>7/8$.
\end{enumerate}

\paragraph{Scope.}
This is a guide to \cite{OAI78}, not an independent proof. We quote the
averages and exponents of \cite{OAI78} as stated there (equation and
line references are to its \TeX\ source), simplify notation, and explain
why each step is needed. We did not re-verify every line of
\cite{OAI78}, which runs to about $16{,}700$ lines of \TeX. Independently, the final statement for
$\zeta$ and for all Dirichlet $L$-functions has a Lean formalization in the
same repository; we compiled it and replayed it through Lean's kernel
(it uses only the standard axioms). The companion guide \cite{guide1112}
explains the simpler $11/12$ argument of \cite{OAI1112}, which is organized
differently.

\section{The family and the number \texorpdfstring{$\beta_*$}{beta*}}\label{sec:family}

Let $\omega=e^{2\pi i/3}$ and let $\OO$ be the Eisenstein integers, with
norm $\N(a+b\omega)=a^2-ab+b^2$. An element is \emph{primary} if it is
$\equiv1\pmod 3$; every ideal prime to $3$ has exactly one primary
generator. For a prime $\pi\nmid6$, the sextic residue symbol $(u/\pi)_6$ is the sixth
root of unity congruent to $u^{(\N(\pi)-1)/6}$ modulo $\pi$ ($0$ if $\pi\mid u$);
for primary $n$ prime to $6$, $(u/n)_6=\prod_{\pi^e\parallel n}(u/\pi)_6^{\,e}$. We write $\chi_n(u)=(u/n)_6$, so that $\chi_n^3$ and
$\chi_n^2$ are the quadratic and cubic symbols.

\paragraph{The family.} A \emph{finite-order Hecke character} $\eta$ of
$\mathbb Q(\sqrt{-3})$ is, concretely, a completely multiplicative function on
primary $n$ that depends only on $n$ modulo some fixed modulus (and is
$0$ when $n$ is not coprime to it); through primary generators it is a function
on ideals. Its $L$-function is
\[
 L(s,\eta)=\sum_{\mathfrak a\ne0}\frac{\eta(\mathfrak a)}{\N(\mathfrak a)^s}
 \qquad(\Re s>1),
\]
summed over the nonzero ideals of $\OO$ (for ideals prime to $3$ this is the sum
over primary $n$; the prime above $3$ enters in the usual way). Examples:
\begin{itemize}
\item $\eta=1$: $L(s,1)=\zeta(s)L(s,\ch)$, the product studied in
 \cite{guide1112};
\item $\eta(n)=\chi(\N(n))$ for a Dirichlet character $\chi$:
 $L(s,\eta)=L(s,\chi)L(s,\chi\ch)$, up to finitely many Euler factors;
\item $\eta(n)=\nu(n)\,\overline{(u/n)_6}$ for fixed $u$ and a
 residue-class character $\nu$: a \emph{sextic twist}. These are the
 $L$-functions that appear as ``rows'' below. In general they are not products of
 Dirichlet $L$-functions.
\end{itemize}
All of these $L$-functions continue meromorphically (entire unless $\eta$ is
principal, when there is a simple pole at $s=1$) and satisfy functional
equations (Hecke). Define
\begin{equation}\label{eq:betastar}
 \beta_*=\sup\bigl(\{1/2\}\cup\{\Re\rho:\ L(\rho,\eta)=0,\
 1/2\le\Re\rho\le1,\ \eta\ \text{primitive of finite order}\}\bigr)
\end{equation}
\cite[L378--385]{OAI78}. Every zero of $\zeta$ is a zero of $L(s,1)$, so
$\beta_*\le7/8$ implies the theorem for $\zeta$.

\paragraph{Why the whole family is unavoidable.} The second evaluation of
the probe produces sextic twists of the target, and the only a priori
information about their zeros is $\Re\rho\le\beta_*$. Moreover, the
contradiction at the end is produced by \emph{whichever} member has a zero
close to $\beta_*$. So even for $\zeta$ every estimate must hold for every
target $\eta$, with margins that do not depend on $\eta$.

\section{The continuation criterion}\label{sec:criterion}

Everything is reduced to the following statement \cite[Prop.\ lem:continuation-criterion]{OAI78}.

\begin{proposition}\label{prop:criterion}
Fix $\sigma_0\in(1/2,1)$, suppose $\beta_*>\sigma_0$, and let
$C(s)=s+c$ with $c$ fixed. Suppose there are $0<\omega<\beta_*-\sigma_0$ and
$\sigma>0$, independent of $\eta$, such that for every primitive $\eta$
there are a finite set $\mathcal S$ of primes, a holomorphic
$H_\eta$ with $|H_\eta(s)-1|\le1/2$ on $\Re s>\sigma_0$, and a function
$J_\eta(Z)$ with
\begin{align}
 |J_\eta(Z)|&\ll_\eta Z^{C(\sigma_0)+\omega}
   &&\text{(low estimate)},\label{eq:low}\\
 |J_\eta(Z)-f_\eta(Z)|&\ll_\eta Z^{C(\beta_*)-\sigma}
   &&\text{(high estimate)},\label{eq:high}
\end{align}
where the \emph{signal} is
\begin{equation}\label{eq:signal}
 f_\eta(Z)=\frac1{2\pi i}\int_{\Re s=2}Z^{C(s)}e^{(s-5/6)^2}
 \frac{H_\eta(s)}{L^{\mathcal S}(s,\eta)}\,ds .
\end{equation}
Then we have a contradiction.
\end{proposition}

\emph{Why.} Together, \eqref{eq:low} and \eqref{eq:high} give
$|f_\eta(Z)|\ll_\eta Z^{C(\beta_*)-\epsilon_*}$ for large $Z$, with
$\epsilon_*=\min(\beta_*-\sigma_0-\omega,\sigma)>0$ independent of $\eta$ and
$\beta_*-\epsilon_*>\sigma_0$. Moving the contour in \eqref{eq:signal} to the right
(the Gaussian decays on vertical lines, and $H_\eta$ and $1/L^{\mathcal S}$ are bounded
on $\Re s\ge2$) gives $f_\eta(Z)\ll_{\eta,B}Z^{B}$ as $Z\to0$, for every $B$.
Hence the Mellin transform $\int_0^\infty f_\eta(Z)Z^{-C(s)}\,dZ/Z$ is holomorphic
in $\Re s>\beta_*-\epsilon_*$, and by Mellin inversion it equals
$e^{(s-5/6)^2}H_\eta(s)/L^{\mathcal S}(s,\eta)$ for $\Re s>1$. Since $|H_\eta|\ge1/2$
on $\Re s>\sigma_0$ and the deleted Euler factors do not vanish in $\Re s>0$,
$1/L(s,\eta)$ continues to $\Re s>\beta_*-\epsilon_*$ for every primitive $\eta$.
But by \eqref{eq:betastar} some primitive $\eta$ has a zero with real part in
$(\beta_*-\epsilon_*,\beta_*]$, a contradiction.

Note the asymmetry: the low estimate is measured against $C(\sigma_0)$ and
the high estimate against $C(\beta_*)$. Since $C$ has slope one,
$C(\beta_*)-C(\sigma_0)=\beta_*-\sigma_0$. The argument in Part~I of
\cite{OAI78} uses $\sigma_0=11/12$ and $C(s)=s-2/3$; the $7/8$ argument
uses
\[
 \sigma_0=\frac78,\qquad C(s)=s-\frac{11}{16},\qquad C(7/8)=\frac3{16}.
\]

\section{The probe}\label{sec:probe}

\subsection{The formula}
Fix the target $\eta$ (for $\zeta$, $\eta=1$). Besides $\eta$, the probe
contains a few fixed factors that depend only on residue classes at a finite set $S$ of
bad primes (the primes above $2$ and $3$, those dividing the modulus of $\eta$, and all
primes of norm below a fixed bound): a character $\xi$ modulo a fixed $b_*$ supported on
$S$, its Gauss sum $\tau$, a reciprocity sign $\mathcal R(\cdot,\cdot)$ and a phase $G$.
The character $\xi$ removes, after Poisson summation, the zero frequency and the
frequencies not prime to $S$; $\tau$, $\mathcal R$ and $G$ cancel unit and reciprocity
phases (between $cn^3$ and $s$, and inside $cn^3$), so that the coefficients after Poisson
summation factor prime by prime. We collect them into
$c_s$ and $\kappa_s(cn^3)$ below (both of absolute value one). With
\[
 \gamma_2(c)=\N(c)^{-1/2}\sum_{x\bmod c}(x/c)_3\,e(x/c),\qquad
 \alpha(c)=c/|c|,\qquad
 e(z)=\exp\bigl(4\pi i\,\Im z/\sqrt3\bigr),
\]
the \emph{base probe} \cite[eq.\ general-probe, L3397]{OAI78} is, after
evaluating its Mellin integral,
\begin{equation}\label{eq:probe}
\boxed{\;
\begin{aligned}
 I_\eta(X,Y,Z)={}&\frac1Y\sum_{s}
 \frac{c_s\,W_1(\N(s)/Y)}{\sqrt{\N(s)X}}
 \sum_{m\in\OO}\xi(m)\,
 \frac{g_s(-m)}{\sqrt{\N(s)}}\,
 W_0\Bigl(\frac{\N(m)}{\N(s)X}\Bigr)\\
 &\times\sum_{\substack{c\ \mathrm{squarefree}\\ n}}
 \frac{\gamma_2(c)\,\overline{\alpha(cn^3)}\,\eta(cn^3)\,\kappa_s(cn^3)\,
       \chi_{cn^3}(m)}{\sqrt{\N(c)}\,\N(n)}\,
 V\Bigl(\frac{\N(c)\N(n)^3}Z\Bigr).
\end{aligned}\;}
\end{equation}
Here $s,c,n$ run over primary elements prime to $S$ (with $s$ not
necessarily squarefree, and $c,n$ allowed to share primes), $m$ runs over all of $\OO$,
\[
 g_s(-m)=\sum_{d\bmod s}(d/s)_6\,e(-md/s)
\]
is a sextic Gauss sum, $W_0,W_1\ge0$ are fixed smooth bumps, and
$V(y)=(2\sqrt\pi)^{-1}e^{-(\log y)^2/4}$ is a log-Gaussian, so that
$\N(cn^3)\approx Z$. (We have absorbed the norm of $b_*$ into $X$; in \eqref{eq:poisson} and in the
principal row below, $X$ is the paper's unabsorbed scale, which changes only bounded factors.)

\paragraph{What the three sums are.}
\begin{itemize}
\item The innermost sum over the \emph{completed index} $cn^3$ is a twisted
 sum of Fourier coefficients of Kubota's cubic theta function, whose
 coefficient at $cn^3$ is (up to normalization) $\gamma_2(c)\N(n)^{1/2}$.
 ``Completed'' means the cube factor $n^3$ is kept; the theta function
 needs it.
\item The sum over $s$, with the Gauss sum $g_s(-m)$, is an additive
 polynomial in the averaging variable $m$.
\item The averaging variable $m$ has $\N(m)\approx XY$; it ties the two
 together through $\chi_{cn^3}(m)$ and $g_s(-m)$.
\end{itemize}
In Part~I of \cite{OAI78}, $X=Y=Z^{1/2}$.

\subsection{First evaluation: the probe is small}\label{sec:low}
Expanding $W_0$ by Mellin inversion separates $m$ from $s$, and splitting $s$ into finitely many
residue classes removes the remaining dependence of $\kappa_s$ on $s$; this gives
an exact bilinear form in the rows $m$ \cite[eq.\ low-separated, L3470]{OAI78}:
\begin{equation}\label{eq:bilinear}
 I_\eta(X,Y,Z)\approx Q^{-1/2}\sum_{\N(m)\asymp Q}A_m B_m,
 \qquad Q=XY,
\end{equation}
(an average over finitely many residue classes and one Mellin variable is
suppressed), where
\[
 A_m=\frac1Y\sum_{\N(s)\asymp Y}c(s)\,\frac{g_s(-m)}{\sqrt{\N(s)}},\qquad
 B_m=\sum_{c,n}\frac{\gamma_2(c)\overline{\alpha(cn^3)}\eta(cn^3)\kappa(cn^3)
       \chi_{cn^3}(m)}{\sqrt{\N(c)}\N(n)}V\Bigl(\frac{\N(cn^3)}Z\Bigr),
\]
with bounded coefficients $c(s)$. Cauchy--Schwarz gives
$|I|\le Q^{-1/2}(\sum|A_m|^2)^{1/2}(\sum|B_m|^2)^{1/2}$, and two large
sieves bound the two mean squares.
\begin{itemize}
\item \emph{Additive large sieve.} $A_m$ is a sum of $e(-md/s)$ over the
 reduced fractions $d/s$, which are $\gg1/Y$ apart in $\mathbb C/\OO$. The
 planar large sieve \cite[Lemma lem:planar-additive-sieve]{OAI78}, with separation
 $\asymp1/Y$ and coefficient mass $O(1/Y)$, gives \cite[Lemma lem:balanced-additive-norm]{OAI78}
 $\sum_{\N(m)\le Q}|A_m|^2\ll(Q+Y^2)/Y$, which is $Q/Y$ when $Y^2\le Q$, i.e.\ $Y\le X$.
\item \emph{Theta function plus quadratic large sieve.} The automorphy of the cubic theta
 function (with Patterson's coefficients as recorded in \cite{DR}) turns $B_m$ into a dual sum
 twisted by a \emph{quadratic} symbol attached to $m$, of length about $\N(m)^2/Z$ (no
 longer than $Z$ when $\N(m)\le Z$), as in \cite[Section 7]{guide1112}; no descent is
 needed, since $B_m$ is already completed. The quadratic large sieve \cite{GL}, with the
 powerful part of $m$ counted separately, gives \cite[Lemma lem:unmarked-completed-moment]{OAI78}
 $\sum_{\N(m)\le Q}|B_m|^2\ll Z^\varepsilon(Q+Q^2/Z)$, i.e.\ the diagonal bound
 $\ll Q\,Z^{\varepsilon}$ when $Q\le Z$.
\end{itemize}
Hence, when $Y\le X$ and $XY\le Z$,
\[
 |I_\eta(X,Y,Z)|\ll Q^{-1/2}\Bigl(\frac QY\Bigr)^{1/2}\bigl(Q\,Z^{\varepsilon}\bigr)^{1/2}
 =X^{1/2}Z^{\varepsilon/2},
\]
and with $X=Y=Z^{1/2}$ this is \cite[Prop.\ prop:balanced-low]{OAI78}
$|I_\eta(Z^{1/2},Z^{1/2},Z)|\ll Z^{1/4+\varepsilon}$.

\subsection{Second evaluation: rows and \texorpdfstring{$1/L(s,\eta)$}{1/L(s,eta)}}\label{sec:high}
Poisson summation in $m$ replaces $m$ by a dual frequency. The character
$\xi$ kills the frequency $0$ and everything not prime to $S$, and every
remaining frequency is uniquely $u\,a^6$ with $u$ a sixth-power-free element,
unit factor included (a \emph{row}), and $a$ primary. After summing over everything else, the
probe becomes \cite[L3831 and eq.\ probe-high]{OAI78}
\begin{equation}\label{eq:poisson}
\boxed{\;
\begin{aligned}
 I_\eta(X,Y,Z)=\frac1{(2\pi i)^3}\iiint
 &X^{1/2-z}Z^{x+z-1}Y^{w-1}e^{(x+z-1)^2}M(z)\widehat W_1(w)\\
 &\times\sixsum_{u}\frac{\overline{\xi(u)}}{\N(u)^{z}}\,
 \frac{\zeta_K^S(6z)\,L^S\bigl(w,\chi_\bullet(u)\bigr)}
      {L^S\bigl(x,\eta\,\overline{\chi_\bullet(u)}\bigr)}\,
 \mathcal H_{\eta,u}(x,w,z)\,dz\,dw\,dx,
\end{aligned}\;}
\end{equation}
where $\chi_\bullet(u)$ is the character $n\mapsto\chi_n(u)=(u/n)_6$,
$\zeta_K=\zeta\cdot L(\cdot,\ch)$, the superscript $S$ removes finitely many
Euler factors, $\mathcal H_{\eta,u}$ is an absolutely convergent Euler
product in explicit regions, and $M(z)=\int_0^\infty\widetilde W_0(r)r^{z-1}\,dr$ and
$\widehat W_1(w)=\int_0^\infty W_1(r)r^{w-1}\,dr$ are Mellin transforms of the weights
($\widetilde W_0$ is the planar Fourier transform of $W_0(|y|^2)$). Each row is therefore a ratio of
$L$-functions: the sum over $a$ gives $\zeta_K(6z)$, the sum over $s$ gives the numerator
$L(w,\chi_\bullet(u))$, and the squarefree variable $c$, whose local factor carries a sign
$-1$ coming from the Gauss sums (as in the M\"obius--Gauss identity of \cite[Section 6]{guide1112}), gives the reciprocal
$1/L(x,\eta\overline{\chi_\bullet(u)})$.

\paragraph{The principal row.} Only $u=1$ has a principal numerator (for a unit
$u\ne1$ the numerator character is a nontrivial Kummer character). Its two simple poles,
at $w=1$ (from $L^S(w,\chi_\bullet(1))=\zeta_K^S(w)$) and at $z=1/6$ (from $\zeta_K^S(6z)$),
give the double residue
\[
 c_S\,X^{1/3}Z^{x-5/6}e^{(x-5/6)^2}\,\frac{H_\eta(x)}{L^S(x,\eta)},\qquad
 H_\eta(x)=\mathcal H_{\eta,1}(x,1,1/6),
\]
with a constant $c_S>0$ \cite[eq.\ shared-principal-data]{OAI78}. After division by $c_S$
and with $X=Z^{1/2}$, this is $Z^{x-2/3}e^{(x-5/6)^2}H_\eta(x)/L^S(x,\eta)$:
exactly the signal \eqref{eq:signal} with $C(s)=C_{\rm I}(s):=s-2/3$. This is where
$C_{\rm I}$ and the Gaussian $e^{(s-5/6)^2}$ come from.

\paragraph{The other rows: zero detection.} A row $u\ne1$ is large only if
its twisted $L$-functions have zeros near the contour. The paper assigns
to each row a bin $a\le\beta_*$: the largest real part of a zero of height at most a
small power $Z^\tau$ among the finitely many twists $\nu\chi_\bullet(u)^{\pm1}$ attached to
the row, rounded down to a fine grid with floor $51/100$. It writes $\delta=2a-1$. A zero $\rho$ of $L(s,\psi)$, $\psi$ the row character,
produces \cite[Prop.\ prop:detector-witness]{OAI78} two large Dirichlet
polynomials with the same character and height: for $U=\N(u)$,
\begin{equation}\label{eq:witness}
 M_r=U^{-r/2}\sum_n\mu(n)\psi(n)W\Bigl(\frac{\N(n)}{U^r}\Bigr),\qquad
 S_m=U^{-m/2}\sum_l\psi(l)W'\Bigl(\frac{\N(l)}{U^m}\Bigr),
\end{equation}
with $t-\tfrac12\le r\le t$, $0\le m\le\tfrac12$ and $r+m\ge t$ for any chosen
$t\in[1,3/2]$, and
\[
 |M_r|^2\gg U^{\delta r-\varepsilon},\qquad |S_m|^2\gg U^{\delta m-\varepsilon}.
\]
(Here $M_r$ is a truncation of $1/L$ and $S_m$ of $L$; at a zero both are
as large as the zero-free region allows.) Part~I uses only $M_r$ with
$t=1$. A large sieve for sextic characters (proved in \cite[Lemma lem:sextic-large-sieve]{OAI78}
by the method of \cite{BGL}) bounds, for weights fixed independently of the row, the mean square
\[
 \sum_{\N(u)\asymp U}\bigl|M_r(u)\bigr|^2
 \ll(UD)^\varepsilon\bigl\{U+D+(UD)^{2/3}\bigr\},\qquad D=U^r,
\]
for squarefree rows (rows with a square-full part are handled by fixing that part and
counting it separately, with the same final result; the dependence of the witness weights on
the row costs only a fixed power of $Z^\tau$). Dividing by the spike $U^{\delta r}$ and taking
the worst length $r\in[1/2,1]$ (the endpoints $r=1/2$ and $r=1$) shows that at most $U^{R(\delta)}$
rows can be bad, with $R(\delta)=\min\{1,\max(1-\delta/2,\,4/3-\delta)\}$
\cite[Prop.\ prop:sextic-row-count]{OAI78}.

\subsection{Why Part I stops at \texorpdfstring{$11/12$}{11/12}}
At the balanced scales $X=Y=Z^{1/2}$ the low estimate is $Z^{1/4}$ and the signal exponent is $C_{\rm I}(\beta_*)=\beta_*-2/3$. The criterion needs
$1/4<\beta_*-2/3$, i.e.\ $\beta_*>11/12$; so
\[
 \frac{11}{12}=\frac23+\frac14 .
\]
At these scales the non-principal rows have room to spare: by our computation from the paper's
formulas, the row bins alone would allow any $\beta_*>5/6$. (The paper's high-side lemmas are,
however, stated only for $\sigma_0\ge7/8$; the Euler region it uses for the principal row and for
$H_\eta$ is $\Re s\ge7/8$.) But the balanced scales are a choice. By Section~\ref{sec:low} the low
estimate is $X^{1/2}$, while the signal carries only $X^{1/3}$, so shrinking $X$ improves the
comparison on the low side. The price is a longer dual sum: rows $u$ up to norm about $Z/X$
(Section~\ref{sec:78probe}). Part~I needs only a bootstrap and does not pursue this trade-off; the
$7/8$ argument is about how far it can be pushed.

\section{The \texorpdfstring{$7/8$}{7/8} probe}\label{sec:78probe}

\subsection{The exponent formula}
With general scales $X=Z^{l_x}$, $Y=Z^{l_y}=Z^{l_x+b}$ and short prime
sums of total log-length $\ell$ (below), the principal row gives
\cite[Def.\ def:stage-high-data]{OAI78}
\[
 C(s)=s-\frac56+\frac{l_x}3+\frac\ell6,
 \qquad h:=1-l_x+\ell,
\]
where $h$ is the effective log-length of the dual rows ($\N(u)\lesssim Z^h$),
and, for the paper's geometry (in particular under the relation $l_x+l_y+\ell=1$ below), the
low estimate becomes $Z^{l_x/2+b/12}$. Equating the two,
\[
 \sigma_0=1-\frac h6+\frac b{12}.
\]
So $\sigma_0$ is governed by the length $h$ of the dual rows: Part~I has $h=1/2$, and $7/8$
needs $h-b/2=3/4$. The paper's design relation $l_x+l_y+\ell=1$, under which its low estimate is
proved, makes this $\sigma_0=11/12-\ell/4$, independent of $b$. Its choices are
\begin{equation}\label{eq:geometry}
 b=\frac18,\quad \ell=\frac16,\quad h=\frac{13}{16},\quad
 l_x=\frac{17}{48},\quad l_y=\frac{23}{48},\qquad
 \sigma_0=\frac{11}{12}-\frac1{24}=\frac78,\quad C(s)=s-\frac{11}{16}.
\end{equation}
\paragraph{Where the gain comes from.} The low side does not need the prime sums. By
Section~\ref{sec:low}, the unmodified probe with $X=Y=Z^{1/4}$ already has low estimate
$Z^{1/8+\varepsilon}$, and the formula above gives $C(s)=s-3/4$; so $\sigma_0=7/8$, with dual
rows of length $h=3/4$. What fails is the high side. There are about $Z^h$ rows, and even the
rows whose $L$-functions have no zeros to the right of $1/2$ must add up to less than the
signal. By our computation from the paper's per-bin exponent
\cite[eq.\ stage-high-exponent]{OAI78}, at real part $1/2$ with every row counted and no prime
factors, these \emph{floor} rows contribute $Z^{C(\sigma_0)+E}$, where
\[
 E=8-9\sigma_0+\frac b6-\ell\qquad\Bigl(\text{when }\sigma_0=1-\frac h6+\frac b{12}\Bigr).
\]
So without slots the paper's row-by-row bounds cannot reach below $\sigma_0=8/9+b/54$, for any
choice of scales; at $X=Y=Z^{1/4}$ they give $E=+1/8$.

The prime slots pay for the longer rows. At fixed $h$ and $b$ they let $X$ and $Y$ both be
$Z^\ell$ larger ($l_x=1+\ell-h$). This leaves $\sigma_0$ unchanged, since the low estimate and
$C(s)$ both rise by $\ell/2$. But it lowers the exponent of every non-principal bin by at least
$\ell$ (our computation from the same formula), because on such a row a slot contributes only a
square-root-normalized prime polynomial (Section~\ref{sec:high78}), while at the principal row
it has its full size $Z^{\ell/6}$. The low side limits the slots: its estimate is proved under
$l_x+l_y+\ell=1$, which together with $\sigma_0=7/8$ forces $\ell=1/6$, whatever $b$ is. With
the paper's geometry the floor rows have $E=-1/48$, and the binding bins are those with zeros
near real part $0.69$ (Section~\ref{sec:final}). Everything else in Part~II exists to make the
high estimate survive rows of length $h=13/16$.

\subsection{Prime slots and the compensated probe}
Fix $K$ disjoint ``slots'' of primes: slot $i$ consists of primes $p\notin S$
in a fixed residue class (the identity class of a fixed finite group)
with $\N(p)\asymp P_i=Z^{\ell_i}$, weighted by a smooth bump $\Omega_i$, where
$\sum\ell_i=\ell=1/6$. Write $I_{\eta;D}$ for the probe \eqref{eq:probe} with the
condition $D\mid cn^3$ inserted in the completed index. The \emph{compensated
probe} is \cite[eq.\ compensated-probe-definition, L6894]{OAI78}
\begin{equation}\label{eq:compensated}
\boxed{\begin{aligned}
 I_{\rm mod}(Z)=\sum_{p_1,\dots,p_K}\prod_{i=1}^K\Omega_i\Bigl(\frac{\N(p_i)}{P_i}\Bigr)
 &\sum_{J\subseteq\{1,\dots,K\}}(-1)^{|J|}\N(p_J)^{-3/2}\,
 \overline{\eta(p_{J^c})}\\
 &\times I_{\eta;\,p_{J^c}}\Bigl(\frac X{\N(p_J)},\frac Y{\N(p_J)},Z\,\N(p_{J^c})\Bigr),
\end{aligned}}
\end{equation}
with $p_J=\prod_{i\in J}p_i$. With one slot it reads
\[
 I_{\rm mod}(Z)=\sum_p\Omega_1\Bigl(\frac{\N(p)}{P_1}\Bigr)\Bigl[
 \underbrace{\overline{\eta(p)}\,I_{\eta;p}(X,Y,Z\N(p))}_{\text{marked}}
 \;-\;\underbrace{\N(p)^{-3/2}I_\eta\bigl(X/\N(p),Y/\N(p),Z\bigr)}_{\text{rescaled}}\Bigr].
\]
The \emph{marked} term restricts the completed index to multiples of $p$
and lengthens its range to $Z\N(p)$. The \emph{rescaled} term shrinks both
averaging scales by $\N(p)$.

\subsection{Low side: still small}
On the low side the marked primes are summed \emph{inside} the completed
row $B_m$, whose mean square stays diagonal, $\ll Q\,Z^\varepsilon$
\cite[Lemma lem:probe-row-norm]{OAI78}; the proof now also uses Heath-Brown's cubic large
sieve \cite{HB} for the marked primes \cite[Lemma lem:hybrid-energy]{OAI78}, and the relation
$l_x+l_y+\ell=1$. The additive mean square acquires a factor from the unequal scales, and each
rescaled slot of length $d$ gains a factor $Z^{-d}$ (from the count $Z^{d}$, the coefficient
$Z^{-3d/2}$ and $X^{1/2}\mapsto X^{1/2}Z^{-d/2}$). The result
\cite[Prop.\ prop:probe-low]{OAI78} is
\[
 |I_{\rm mod}(Z)|\ll Z^{l_x/2+b/12+\varepsilon}=Z^{17/96+1/96+\varepsilon}=Z^{3/16+\varepsilon}.
\]
The term $b/12$ is the price of the unequal scales. With $P_a=Y^2/Q=Y/X\approx Z^{b}=Z^{1/8}$,
the paper's Gram bound gives $(Q/Y)\bigl(1+P_a^{1/6}+P_a^2/Y\bigr)$ for the additive mean square
\cite[Prop.\ prop:probe-gram]{OAI78} (the planar sieve alone would give $(Q/Y)(1+P_a)$); the
middle term counts the rare frequencies that are sixth powers up to factors supported on the
common divisor and on $S$. Since $P_a^2/Y\ll P_a^{1/6}$, the middle term dominates, and after
Cauchy--Schwarz it becomes $P_a^{1/12}=Z^{b/12}$.

\subsection{High side: what the compensation cancels}\label{sec:high78}
In the Poisson evaluation \eqref{eq:poisson}, a slot prime $p\nmid u$
multiplies row $u$ by (roughly)
\[
 \text{marked: }\ -\overline{\chi_p(u)}+\N(p)^{-w},\qquad
 \text{rescaled: }\ \N(p)^{-w},
\]
(times a common factor $\N(p)^{z-1}$). The difference is
$-\overline{\chi_p(u)}$ plus small terms
\cite[eq.\ old-eq:5.13b, Prop.\ prop:probe-errors]{OAI78}. (For the few slot primes
$p\mid u$ the main term vanishes and $-\N(p)^{-w}$ survives; the paper controls them through the
drop in the row's conductor \cite[eqs.\ old-eq:5.13d, old-eq:5.13e]{OAI78}.) So each slot
contributes, in row $u$, the factor $-P_i^{z-1/2}Q_i(u)$, where
\begin{equation}\label{eq:amplitude}
 Q_i(u)=P_i^{-1/2}\sum_{p\in\text{slot }i}\Omega_i\Bigl(\frac{\N(p)}{P_i}\Bigr)
 \Bigl(\frac{\N(p)}{P_i}\Bigr)^{z-1}\overline{\Bigl(\frac up\Bigr)_6}
\end{equation}
is a centrally normalized \emph{prime polynomial}, typically of size $O(1)$. Without the subtraction, the
row-independent term $\N(p)^{-w}$ would act on \emph{every} row like a
prime polynomial of maximal size, and could not be controlled by any
average over rows.

At the principal row $u=1$ each slot gives $-1$ per prime. Hence the
signal survives, multiplied by the nonzero normalizer
\[
 A_T(Z)=(-1)^KZ^{-\ell/6}\prod_{i=1}^K\sum_{p\in\text{slot }i}
 \Omega_i\Bigl(\frac{\N(p)}{P_i}\Bigr)\N(p)^{-5/6},
\]
with $|A_T(Z)|\asymp(\log Z)^{-K}$, so that dividing by it costs only $Z^{\varepsilon}$
($A_T(Z)>0$ when $K$ is even), and the normalized probe
$J_{\rm II,\eta}=I_{\rm mod}/(c_SA_T(Z))$ has signal \eqref{eq:signal} with
$C(s)=s-11/16$; the factor $Z^{\ell/6}$ in $C$ is the slots evaluated at $z=1/6$.

\section{Counting the bad rows}\label{sec:counts}

\subsection{Three averages}
A row in bin $\delta$ has prime polynomials of \emph{mean amplitude} $q$: the length-weighted
mean of the exponents $g_i$ with $|Q_i|\approx P_i^{g_i}$. The bin ceiling
$|Q_i|\ll P_i^{\delta/2+o(1)}$ gives $0\le q\le\delta/2$, and suitably selected slots of total
base-$U$ length $\lambda$ satisfy $|\prod_{\rm sel}Q_i|^2\ge U^{2q\lambda-o(1)}$. Such a row
is detected by large values of $M_r$, $S_m$ and the $Q_i$. Three
averages over rows bound how many such rows exist.

\paragraph{1. Inverse moment with prime factors}
\cite[Lemma lem:marked]{OAI78}. With row range $H$, length $D$, slot sizes $P_i$, $P=\prod_iP_i$, and fixed
residue-class characters $\nu$ and coefficients $a_i(p)$,
\begin{equation}\label{eq:inverse-moment}
 \sum_{0<\N(u)\ll H}\Bigl|D^{-1/2}\sum_n\mu(n)\nu(n)\Bigl(\frac un\Bigr)_6W\Bigl(\frac{\N(n)}D\Bigr)
 \prod_iP_i^{-1/2}\sum_{p\in\mathcal P_i}a_i(p)\nu(p)\Bigl(\frac up\Bigr)_6\Omega_i\Bigl(\frac{\N(p)}{P_i}\Bigr)\Bigr|^2
 \ll H^{1+\varepsilon},
\end{equation}
provided $DP^2\le HZ^{-c_1}$ and $D^2P^8\le H^3Z^{-c_2}$ (fixed $c_1,c_2>0$). With no prime factors this
is exactly the mean-square estimate behind the alternate $11/12$ proof of \cite{OAI1112}
(explained in \cite{guide1112}; Part~I of \cite{OAI78} does not use it),
$\sum_{\N(u)\le D^{1+\vartheta}}|\sum_n\mu(n)(u/n)_6W(\N(n)/D)|^2\ll D^{2+\vartheta+\varepsilon}$,
and its proof follows the same route: Poisson in $u$, the identity converting
M\"obius coefficients into cubic Gauss sums, the cubic theta function with
the quadratic large sieve, and the two-Poisson descent \cite[Sections 6--8]{guide1112}. With
prime factors the paper adds a reduction for slot primes dividing $n$ and, in the terminal step,
Heath-Brown's cubic large sieve for the slot primes \cite[Lemma lem:hybrid-energy]{OAI78}.

\paragraph{2. Plain fourth moment with prime factors}
\cite[Lemma lem:plain]{OAI78}. With lengths $Z^{n_1},Z^{n_2}$ and rows
$\N(k)\ll Z^{m'}$ (excluding rows whose character is principal, or lies in a fixed finite group, when prime factors are present),
\begin{equation}\label{eq:plain-moment}
 \sum_{k}\Bigl|\prod_{j=1}^2Z^{-n_j/2}\sum_l\tau(l)\Bigl(\frac kl\Bigr)_6W_j\Bigl(\frac{\N(l)}{Z^{n_j}}\Bigr)
 \cdot\prod_iP_i^{-1/2}\sum_p\tau(p)\nu_i(p)\Bigl(\frac kp\Bigr)_6\Omega_i\Bigl(\frac{\N(p)}{P_i}\Bigr)\Bigr|^2
 \ll Z^{m'+\varepsilon},
\end{equation}
for all lengths when there are no prime factors, and otherwise, for $3/4\le\kappa\le1$ and slots
of sufficiently small individual length, when
$n_1+n_2+6\kappa\lambda\le m'$, where $\kappa=2\beta_*-1$ (we state the case used in the application, where
the paper's moving-support parameter is $0$). This is new. It is an induction on the effective width: the functional equation of the
sextic twists $L(s,\psi)$ (Hecke reflection) shortens the sums, and two finite Poisson
transforms per step return products of the same kind at smaller width. Transformed rows whose
character lies in the fixed finite group can carry volume-sized main terms, which an
equal-scale centered difference cancels. The zero-free hypothesis $\beta_*\le(1+\kappa)/2$
(which holds, with equality, for $\kappa=2\beta_*-1$) bounds the prime factors pointwise,
$|\prod_iQ_i|^2\ll Z^{\kappa\lambda+\varepsilon}$.

\paragraph{3. Sixth-power amplification}
\cite[Lemma lem:inverse-amplification]{OAI78}. Replacing a sixth-power-free row $u$ ($\N(u)\asymp U$) by
$ua^6$ for the $\asymp(H/U)^{1/6}$ ideals $a$ with $\N(a)\le(H/U)^{1/6}$, where
$H=\max(2U,D^{1+c})$, and applying (1) without prime factors at row range $H$ (the same trick
as in the deduction of $11/12$ in \cite{OAI1112,guide1112}), gives
$\sum_{\N(u)\asymp U}|M_r(u)|^2\ll U^{e(r)+\varepsilon}$ with $D=U^r$ and
$e(r)=\max\{1,(1+5r)/6\}$.

\subsection{Combining the counts}
Apply the moments at base $U$: rows have $\N(u)\asymp U$, so the row exponent is $1$, and $r$,
$m$, $\lambda$ are base-$U$ lengths. Dividing each moment by the guaranteed size of a bad row
($|M_r|^2\gg U^{\delta r}$, $|S_m|^4\gg U^{2\delta m}$, $|\prod_{\rm sel}Q_i|^2\gg U^{2q\lambda}$)
gives, for $\#\mathcal B$ rows of norm $\asymp U$ in one bin,
\[
 \#\mathcal B\ll U^{1-\delta r-2q\lambda+\varepsilon}\quad\text{(inverse)},\qquad
 \#\mathcal B\ll U^{1-2\delta m-2q\lambda+\varepsilon}\quad\text{(plain, two copies of }S_m),
\]
with the prime-factor ``capacities'' $\lambda\le(1-r)/2$ (from $DP^2<H$; the second condition
then needs $r>1/2$, and the inverse count is used only for $r\ge23/37$) and
$\lambda\le(1-2m)/(6\kappa)$ (from $2m+6\kappa\lambda\le1$).
The inverse bound is good when the M\"obius witness is long, the plain bound when the plain witness is long;
since $r+m\ge t$, at least one witness is long, and the better of the two bounds is used. Writing $q=x\delta$ with
$0\le x\le1/2$, taking the better of the two, and using the amplified count
$L(t)=1-\delta+(\alpha-\delta)(t-1)$, $\alpha=5/6$, for long inverse witnesses, gives
\cite[Prop.\ prop:detector-counts]{OAI78} (assuming the available prime length exceeds $7/37$ by a fixed
margin; at the longest rows it is $\ell/h=8/39$)
\[
 \#\mathcal B\ll U^{\max\{R_{\rm short}(t),L(t)\}+\Delta/4+\varepsilon},\qquad
 R_{\rm short}(t)=1-\delta+\frac{\delta P_x}{D_x}\Bigl(\frac32-t\Bigr),
\]
where $D_x=3-17x/9$, $P_x=(2-8x/9)(1-x)$, and $\Delta=\beta_*-7/8$.

\section{The final inequality}\label{sec:final}

A bin of rows of norm $U=Z^d$ contributes at most $Z^{C(7/8)+E(d)+o(1)}$ to the
probe, where \cite[eq.\ common-high-exponent]{OAI78}
\[
 E(d)=-\frac1{48}+\frac23\delta+\frac q6-h(1-R)
 +(d-h)\Bigl(R+\frac\delta2-\frac{17}{50}\Bigr).
\]
Its terms record the Mellin weight on the central contour, the size
$U^{\delta/2}$ of the numerator $L(w,\chi_\bullet(u))$ on the reflected line,
the number $U^R$ of rows, and the prime polynomials. The slope in $d$ is
positive, so the worst rows are the longest, $d=h=13/16$. There one takes
$t=1+\delta P_x/(2\mathcal J)$ with
$\mathcal J=(\alpha-\delta)D_x+\delta P_x$, so that $R=R_*=1-\delta+(\alpha-\delta)\delta P_x/(2\mathcal J)$, and the
paper proves \cite[Lemma lem:balanced-endpoint]{OAI78}
\[
 -E(h)\ \ge\ \frac{49}{440640}\qquad\text{for all }0\le\delta\le\frac56,\ 0\le x\le\frac12 .
\]
Including the small loss from $\kappa=2\beta_*-1$, the comparison actually needed is
\[
 E_{\rm actual}(h)-\Delta\le-\frac{49}{440640}-\frac{51}{64}\Delta+\epsilon_{\rm total}<0,
\]
where $\epsilon_{\rm total}$ collects the adjustable losses (chosen small after the certificate);
i.e.\ every non-principal bin is smaller than the signal size $Z^{C(\beta_*)}=Z^{C(7/8)+\Delta}$ by a fixed power.
Rows slightly longer than $Z^h$ ($h<d\le h+\zeta$) cost at most $2\zeta$, absorbed by taking
$\zeta$ small. The top bin $\delta=5/6$ (no selected primes; $E(h)\le-1/48-\delta/16$) and the
floor bin ($E\le-7/1200$) are separate easy cases, and shorter rows ($d\le1/2$, no prime factors
needed), very short rows and very long rows ($d>h+\zeta$) are handled with room to spare. All margins are chosen before the target $\eta$, so they are uniform in $\eta$, and
Proposition~\ref{prop:criterion} with $\sigma_0=7/8$ gives the contradiction. Hence $\beta_*\le7/8$, and
$\zeta(s)\ne0$ for $\Re s>7/8$.

\paragraph{The worst case.} At $x=1/2$ (prime polynomials as large as the bin allows) and
$\delta\approx0.387$ (zeros with real part about $0.69$) one finds
$R_*\approx0.668$ and $E(h)\approx-2.3\times10^{-4}$: the inequality holds, but only just.

\section{What each ingredient does}\label{sec:summary}
\begin{center}
\begin{tabular}{p{4.2cm}p{10.3cm}}
\toprule
Ingredient & Role\\\midrule
Family supremum $\beta_*$ & Gives a priori control of every twisted $L$-function that appears as a row; the contradiction may come from any member.\\
Bootstrap $\beta_*\le11/12$ & From Part~I. Gives $\Delta\le1/24$ and $\kappa=2\beta_*-1\le5/6$; since $\delta\le2\beta_*-1$ by the definition of $\beta_*$, $\delta\le5/6$ for every bin, as used in the counts and the final certificate.\\
Prime slots, $\ell=1/6$ & Make the long dual rows ($h=13/16$) affordable: at fixed $h$ they let $X,Y$ be $Z^\ell$ larger, which leaves $\sigma_0=1-h/6+b/12$ unchanged and lowers every non-principal bin; full size at $u=1$, square-root-size prime polynomials elsewhere. The low side caps them at $\ell=1/6$.\\
Compensation (marked $-$ rescaled) & Removes the row-independent term $\N(p)^{-w}$, leaving the row polynomial $-\overline{\chi_p(u)}$; free on the low side; keeps the principal row nonzero.\\
Unequal scales, $b=1/8$ & Cost $Z^{b/12}$ on the low side, gain on the high side; $b$ cancels from $\sigma_0$ under $l_x+l_y+\ell=1$.\\
Two witnesses $M_r,S_m$ & Complementary detectors from one zero ($r+m\ge t$): a long M\"obius witness is counted by the inverse moment (or by sixth-power amplification when $r\ge1$), a long plain witness by the plain fourth moment.\\
Inverse moment & The machinery of the alternate $11/12$ proof \cite{OAI1112} (cubic theta, quadratic large sieve, two-Poisson descent), extended to prime factors with Heath-Brown's cubic large sieve.\\
Plain fourth moment & New: an induction combining Hecke reflection with two finite Poisson transforms (and a centered difference for rows in the fixed finite group); prime factors bounded pointwise using $\beta_*\le(1+\kappa)/2$, which holds with equality for $\kappa=2\beta_*-1$.\\
\bottomrule
\end{tabular}
\end{center}

\paragraph{How tight is it?} The certificate guarantees a margin of $49/440640\approx1.1\times10^{-4}$;
the true worst case is about $2.3\times10^{-4}$. We re-evaluated the paper's general exponent \cite[eq.\ stage-high-exponent]{OAI78} at $d=h$
with the same row counts (which do not depend on $\ell$ or $b$ while $\ell/h>7/37$), together with
$\sigma_0=1-h/6+b/12$ and $l_x+l_y+\ell=1$ (our computation, not a claim of \cite{OAI78}). The
endpoint inequality then survives only for $\ell\lesssim0.1668$ when $b=1/8$, and only for
$0.094\lesssim b\lesssim0.152$ when $\ell=1/6$; at $\ell=0.17$ or $b=0$ the worst bin exceeds the
signal by about $Z^{0.004}$. So $7/8$ appears to be close to the limit of this particular design.

\begin{thebibliography}{9}
\bibitem{OAI78} OpenAI, \emph{The Quasi-Riemann Hypothesis: A Zero-Free Half-Plane $\Re s>7/8$},
manuscript dated September 30, 2026,
\href{https://github.com/openai/math/tree/main/preprints/The-Quasi-Riemann-Hypothesis-September-30-2026}{github.com/openai/math}.
Line references are to \texttt{build/paper.tex}.
\bibitem{OAI1112} OpenAI, \emph{The Quasi-Riemann Hypothesis}, manuscript dated October 5, 2026
(an alternate proof of the $11/12$ half-plane), same repository.
\bibitem{guide1112} Claude Opus 5.5 and OpenAI Astra, \emph{Sixth powers in a mean square: an accessible
guide to the $11/12$ argument for $\zeta(s)$}, companion note,
\url{https://www.math.uwaterloo.ca/~mrubinst/quasi_rh_notes/zeta_11_12_guide.pdf}.
\bibitem{BGL} V.~Blomer, L.~Goldmakher and B.~Louvel, \emph{$L$-functions with $n$-th-order twists},
Int.\ Math.\ Res.\ Not.\ IMRN (2014), no.~7, 1925--1955.
\bibitem{DR} A.~Dunn and M.~Radziwi\l\l, \emph{Bias in cubic Gauss sums: Patterson's conjecture},
Ann.\ of Math.\ (2) \textbf{200} (2024), 967--1057.
\bibitem{GL} L.~Goldmakher and B.~Louvel, \emph{A quadratic large sieve inequality over number fields},
Math.\ Proc.\ Cambridge Philos.\ Soc.\ \textbf{154} (2013), 193--212.
\bibitem{HB} D.~R.~Heath-Brown, \emph{Kummer's conjecture for cubic Gauss sums}, Israel J.\ Math.\
\textbf{120} (2000), 97--124.
\end{thebibliography}
\end{document}
